
\documentclass[11pt]{article}
\usepackage[utf8]{inputenc}
\usepackage[T1]{fontenc}
\usepackage[a4paper,margin=1in]{geometry}
\usepackage{lmodern}
\usepackage{microtype}
\usepackage{booktabs}
\usepackage{longtable}
\usepackage{tabularx}
\usepackage{array}
\usepackage{enumitem}
\usepackage{hyperref}
\usepackage{xcolor}
\usepackage{amsmath}
\usepackage{float}
\hypersetup{colorlinks=true,linkcolor=blue,citecolor=blue,urlcolor=blue}

\title{\textbf{Speech Recognition Deep Learning Project Plan}\\
\large TensorFlow Speech Recognition Challenge / Speech Commands v0.01}
\author{}
\date{\today}

\begin{document}
\maketitle


\section{Dataset and task definition}

The training archive contains 1-second, 16 kHz spoken-word clips from 30 word folders -- \texttt{yes, no, up, down, left, right, on, off, stop, go, bed, bird, cat, dog, eight, five, four, happy, house, marvin, nine, one, seven, sheila, six, three, tree, two, wow, zero}. The evaluation target is only the 10 commands, plus the special classes \texttt{unknown} and \texttt{silence}. The \texttt{\_background\_noise\_} folder contains long ambient recordings that should be randomly cut into 1-second segments to form silence training examples. The project will therefore study a 12-class classification task:
\begin{itemize}[leftmargin=*]
    \item \textbf{Target commands}: \texttt{yes, no, up, down, left, right, on, off, stop, go}
    \item \textbf{Special classes}: \texttt{unknown}, \texttt{silence}
\end{itemize}


For research experiments, the project should use the official \texttt{validation\_list.txt} and \texttt{testing\_list.txt} split files. This follows the reproducible protocol proposed for Speech Commands and avoids speaker leakage between train/validation/test.


\section{Network architectures}

The project compare three model families:

\begin{enumerate}[leftmargin=*]
    \item \textbf{KWT-1} - Transformer based model.
    \item \textbf{BC-ResNet-1} - compact CNN and the main convolutional model.
    \item \textbf{MatchboxNet-3x1x64} - efficient 1D time-channel separable CNN.
\end{enumerate}

These models are modern enough to give competitive results on Speech Commands while remaining small enough for repeated experimentation \cite{bcresnet,kwt,matchbox}. A summary is given in Table \ref{tab:models}.

\begin{table}[H]
\centering
\small
\begin{tabularx}{\textwidth}{>{\raggedright\arraybackslash}p{2.5cm} >{\raggedright\arraybackslash}p{2.4cm} >{\raggedright\arraybackslash}p{1.8cm} >{\raggedright\arraybackslash}X}
\toprule
Model & Default paper frontend & Parameters & Default training setup used in the literature \\
\midrule
BC-ResNet-1 \cite{bcresnet} &
40-dim log-Mel, 30 ms window, 10 ms hop &
9.2k &
200 epochs; SGD with momentum 0.9; weight decay \(10^{-3}\); batch size 100; LR warmup to 0.1 for 5 epochs then cosine decay; dropout 0.1; time shift \([-100,100]\) ms; background-noise mixing; SpecAugment. \\
\addlinespace
KWT-1 \cite{kwt,kwtrepo} &
40-dim log-Mel, 30 ms window, 10 ms hop &
607k &
23,000 steps; AdamW; LR \(10^{-3}\); batch size 512; cosine schedule; 10 warmup epochs; weight decay 0.1; label smoothing 0.1; time shift, resampling, background perturbation and SpecAugment. \\
\addlinespace
MatchboxNet-3x1x64 \cite{matchbox} &
64 MFCC, 25 ms windows, 10 ms overlap, padded to 128 frames &
77k &
200 epochs; NovoGrad (\(\beta_1=0.95, \beta_2=0.5\)); max LR 0.05, min LR 0.001; warmup-hold-decay schedule; weight decay \(10^{-3}\); mixed precision; batch size 128 per GPU; time shift, white noise, SpecAugment and SpecCutout. \\
\bottomrule
\end{tabularx}
\caption{Models and their default paper settings}
\label{tab:models}
\end{table}



\section{Audio preprocessing}

\subsection{Primary models frontend}

To isolate architectural differences, the \textbf{primary comparison should use one shared frontend}: 40-bin log-Mel spectrograms extracted from 1-second clips using a 30 ms window and 10 ms hop, which yields approximately 98 time frames by 40 frequency bins per example. Each clip should be padded or cropped to 1 second before feature extraction.

\subsection{Baseline model frontend}

MatchboxNet was originally reported with a different frontend: 64 MFCC features computed from 25 ms windows with 10 ms overlap, padded to 128 feature vectors per sample \cite{matchbox}.

\subsection{Augmentation}

A good shared augmentation baseline is:
\begin{itemize}[leftmargin=*]
    \item random time shift,
    \item random mixing with background-noise clips from \texttt{\_background\_noise\_},
    \item SpecAugment (time masking, frequency masking)
\end{itemize}

BC-ResNet and KWT both use time shifting, background perturbation and SpecAugment-like masking; MatchboxNet also uses time shift plus strong spectrogram masking and cutout \cite{bcresnet,kwt,matchbox}.

\section{Handling the \texttt{unknown} classes}

This is the most important methodological part of the project. The standard literature setup is to evaluate on 12 classes, where the remaining 20 training words are treated as \texttt{unknown} and silence is synthesized from background noise. The project compare four strategies.

\subsection*{Strategy A - post-hoc remapping}

Train on all 31 (30 commands + \texttt{silence}) original labels, then remap any non-target prediction to \texttt{unknown} at inference time. This strategy is simple, but it is not well aligned with the final objective because the model spends capacity separating the 20 non-target words even though the final task only asks whether they belong to a single \texttt{unknown} bucket.

\subsection*{Strategy B - collapsed training labels (recommended main baseline)}

Collapse all 20 non-target words into a single \texttt{unknown} label before training, add \texttt{silence} examples from \texttt{\_background\_noise\_}, and train a 12-class classifier directly. This is the most standard and task-aligned protocol, and it is the setup followed by BC-ResNet and many Speech Commands papers \cite{bcresnet}. 

To prevent the \texttt{unknown} category from dominating training, non-target word classes will be rebalanced so that the total number of \texttt{unknown}/outlier examples is approximately equal to the average size of the target keyword classes. Examples will be sampled in a class-aware manner, drawing a similar number of samples from each non-target word class to preserve diversity within the \texttt{unknown} category and avoid over-representation of a few frequent classes.

\subsection*{Strategy C - two-stage classification}

A hierarchical two-head architecture is used. A shared backbone feeds two classification heads. Head~1 performs 3-way classification into \texttt{silence}, \texttt{unknown}, and \texttt{keyword}. Head~2 performs 10-way classification over the target keywords. During training, the loss of Head~1 is computed for all samples using standard cross-entropy with the corresponding 3-class target. The loss of Head~2 is computed only for samples whose ground-truth label belongs to the 10 target keywords; for \texttt{silence} and \texttt{unknown} samples, the Head~2 loss is masked out and does not contribute to optimization. The total loss is therefore
\[
\mathcal{L} = \mathcal{L}_{\text{head1}} + \lambda \,\mathbf{1}[y \in \mathcal{K}]\, \mathcal{L}_{\text{head2}},
\]
where \(\mathcal{K}\) denotes the set of 10 target keywords, and \(\lambda\) is a scalar weight controlling the contribution of the second head (\(\lambda\) = 1 at first, and tuned later if needed). During inference, the prediction of Head~1 is used first; Head~2 is consulted only when Head~1 predicts \texttt{keyword}.

For the hierarchical 3-way head (\texttt{silence}/\texttt{unknown}/\texttt{keyword}), class imbalance will be handled by balanced sampling minibatches so that silence, unknown, and keyword are sampled in roughly equal proportions.

\subsection*{Strategy D - confidence-based rejection with uniform targets}

This strategy is inspired by \textit{Deep Anomaly Detection with Outlier Exposure} \cite{oe}.
\\

Train only an 11-class classifier for the known classes (\texttt{10 commands + silence}). For non-target training words, encourage a uniform output distribution over the 11 known classes (i.e., the target label for unknown class becomes \([1/11, \dots, 1/11]\)). Then, classify a sample as \texttt{unknown} at inference time if
\[
\max_k p(y = k \mid x) < \tau.
\]
The threshold \(\tau\) is tuned on the validation set. For fairness, the non-target word pool will use the same class-aware balancing policy as in Strategy~B, so that outlier exposure remains diverse without overwhelming the target keyword classes.


\section{Experimental protocol}

The study is organized in two stages. In the first stage, each of the three model families will be trained from scratch under all four unknown-handling strategies using literature-based default settings, the official \texttt{train}/\texttt{validation}/\texttt{test} split files, and each model's standard frontend. BC-ResNet-1 and KWT-1 will use 40-bin log-Mel features, while MatchboxNet-3x1x64 will use its native 64-dimensional MFCC frontend. For Strategy~D, the confidence threshold \(\tau\) will be tuned on the validation set, and a threshold sweep will be reported to show how the validation metrics vary with \(\tau\). In this stage, all models will be trained with a minimal shared augmentation baseline consisting of random time shift and background-noise mixing, while SpecAugment is disabled.

In the second stage, augmentation will be studied only for the best unknown-handling strategy of each backbone using a \(2\times2\times2\) grid (Table~\ref{tab:aug_grid}).

\begin{table}[h]
\centering
\caption{Augmentation grid factors and their off/on settings.}
\label{tab:aug_grid}
\begin{tabular}{llp{6.5cm}}
\toprule
\textbf{Factor} & \textbf{Off} & \textbf{On} \\
\midrule
Time shift & --- & Uniform random shift in $[-100, 100]$\,ms \\
Background noise & --- & Mixed with $p{=}0.8$; random 1\,s crop from \texttt{\_background\_noise\_}, gain $\alpha \sim U(0, 0.3)$ \\
SpecAugment & --- & One time mask (max 10 frames) $+$ one frequency mask (max 5 bins) \\
\bottomrule
\end{tabular}
\end{table}

As an additional follow-up experiment, the best-performing KWT Stage~2 setting will be reused and varying only the KWT hyperparameters \(\texttt{attention\_type} \in \{\texttt{time}, \texttt{freq}, \texttt{both}\}\), \(\texttt{depth} = 6\), and \(\texttt{num\_heads} \in \{1, 2\}\). The goal of this extra study is to test alternative KWT hyperparameter settings while reducing the total number of parameters by roughly \(2\times\) for the reduced-depth \texttt{time} and \texttt{freq} variants.

Model selection will be based primarily on validation 12-class macro-F1. Final test-set results will report: 12-class accuracy; 12-class macro-F1; macro-F1 over the ten target commands only; per-class precision, recall, and F1; and raw and row-normalized confusion matrices for all 12 classes.



\begin{thebibliography}{99}

\bibitem{bcresnet}
B. Kim, M. Lee, J. Lee, S. Yang, and K. Hwang, ``Broadcasted Residual Learning for Efficient Keyword Spotting,'' Interspeech, 2021. \url{https://arxiv.org/abs/2106.04140}

\bibitem{kwt}
A. Berg, M. O'Connor, and M. Tairum Cruz, ``Keyword Transformer: A Self-Attention Model for Keyword Spotting,'' Interspeech, 2021. \url{https://arxiv.org/abs/2104.00769}

\bibitem{kwtrepo}
ARM Software, ``keyword-transformer'' (official implementation and training scripts). \url{https://github.com/ARM-software/keyword-transformer}

\bibitem{matchbox}
S. Majumdar and B. Ginsburg, ``MatchboxNet: 1D Time-Channel Separable Convolutional Neural Network Architecture for Speech Commands Recognition,'' Interspeech, 2020. \url{https://arxiv.org/abs/2004.08531}

\bibitem{oe}
D. Hendrycks, M. Mazeika, and T. Dietterich, ``Deep Anomaly Detection with Outlier Exposure,'' arXiv:1812.04606, 2018. \url{https://arxiv.org/abs/1812.04606}

\end{thebibliography}

\end{document}
